\documentclass[a4paper,10pt]{amsart}

\pdfinfoomitdate=1
\pdftrailerid{}
\pdfsuppressptexinfo=15

\usepackage[T1]{fontenc}
\usepackage{lmodern}
\usepackage[margin=25mm]{geometry}
\usepackage{amsmath,amssymb,mathtools,booktabs,array,xcolor}
\usepackage[protrusion=true,expansion=false]{microtype}
\usepackage[shortlabels]{enumitem}
\usepackage[numbers,sort&compress]{natbib}
\usepackage[colorlinks=true,linkcolor=blue!55!black,citecolor=blue!55!black,urlcolor=blue!55!black]{hyperref}
\hypersetup{hypertexnames=false,pdftitle={},pdfauthor={},pdfsubject={},pdfkeywords={},pdfcreator={},pdfproducer={}}

\newtheorem{theorem}{Theorem}
\newtheorem{lemma}[theorem]{Lemma}
\newtheorem{proposition}[theorem]{Proposition}
\newtheorem{corollary}[theorem]{Corollary}
\theoremstyle{remark}
\newtheorem{remark}[theorem]{Remark}

\newcommand{\F}{\mathbb F}
\newcommand{\K}{K}
\newcommand{\X}{\mathcal X}
\newcommand{\J}{\mathcal J}
\newcommand{\Rec}{\operatorname{Rec}}
\newcommand{\Fix}{\operatorname{Fix}}
\newcommand{\Span}{\operatorname{span}}
\newcommand{\im}{\operatorname{im}}
\newcommand{\tail}{\operatorname{tail}}
\newcommand{\Tr}{\operatorname{Tr}}

\title[Quartic inverse-span dynamics]{Quartic Inverse-Span Dynamics:\ Binary Depth Jump and Exact Fibres}
\author{Anonymous}
\date{}

\begin{document}

\begin{abstract}
Let $p$ be prime and send each $\F_p$-subspace $A$ of
$\F_{p^4}$ to the span of the inverses of its nonzero points.  Published
work already classifies when patched inversion carries a subspace onto a
subspace and describes inverse projective lines; those facts receive no
contribution credit here.  We use them to determine the resulting finite
functional graph.  Its sharp maximum tail is two at $p=2$ and one at every
odd prime: a non-subfield plane passes through a hyperplane only in the
binary case.  The recurrent states are zero, the full field, all lines, and
all scalar copies of the quadratic subfield, and every period divides two.
We give the exact depth enumerator, image stabilization, fixed-iterate counts,
and dynamical zeta function.  We also determine every target's fibre at every
positive time.  In particular, the $30$ transient planes over $\F_2$
distribute uniformly over the $15$ hyperplanes, two predecessors per target,
before reaching the full field.  A standalone exhaustive control checks
$32{,}754$ assertions for $p=2,3,5$.  This owner-dependent Round-0 note
remains on external hold.
\end{abstract}

\maketitle

\section{The literal map, published inputs, and theorem ceiling}

Fix a prime $p$, set $\K=\F_{p^4}$, and let $\X_p$ be the lattice of all
$\F_p$-linear subspaces of $\K$.  For $A\in\X_p$, define
\begin{equation}\label{eq:map}
 \J(A)=\Span_{\F_p}\{a^{-1}:a\in A\setminus\{0\}\},
 \qquad \J(0)=0.
\end{equation}
The span in \eqref{eq:map} is part of the update.  Powers $\J^t$ mean
iteration on the subspace lattice, not multiplicative powers of its points.

Two published inputs set a hard ceiling on the claims.  First,
Kolomeec and Bykov classify the affine $\F_p$-subspaces whose image under
patched inversion $0\mapsto0$ is again affine: apart from the size-two
boundary, they are exactly scalar copies of subfields
\citep{KolomeecBykov2024}.  Earlier inverse-closed and large-intersection
results provide the algebraic context \citep{Mattarei2007,Csajbok2013}.
Second, inverse projective lines in the cyclic field model are known to be
normal rational curves in their spans; the small-field regime gives
independent tuples \citep{FainaEtAl2002,LavrauwZanella2014}.  We use the
classification as an external proposition and include a short
denominator-clearing calculation for the line span.  Neither input is a
contribution of this note.

For $A\in\X_p$, let $\tail(A)$ be the least time at which its orbit enters a
cycle.  Put
\begin{align}\label{eq:constants}
 L&=p^3+p^2+p+1, &
 P&=(p^2+1)(p^2+p+1), & Q&=p^2+1,\\
 S&=2+2L+P, & R&=2+L+Q, &
 F&=2+\gcd(2,L)+\gcd(2,Q). \notag
\end{align}
Thus $L$ is the common number of lines and hyperplanes, $P$ is the number
of planes, and $S=|\X_p|$.
For $t\geq1$ and $B\in\X_p$, write
\[
 \nu_t(B)=|\{A\in\X_p:\J^t(A)=B\}|.
\]
The finite-map zeta function is understood formally as
\[
 \zeta_{\J}(z)=\exp\!\left(\sum_{n\geq1}
          |\Fix(\J^n)|\frac{z^n}{n}\right),
\]
following the periodic-point convention of \citet{ArtinMazur1965}.

\begin{theorem}[complete quartic dynamics]\label{thm:main}
For the map \eqref{eq:map}, the following statements hold.
\begin{enumerate}[(i)]
\item The recurrent states are exactly
\begin{equation}\label{eq:recurrent}
 0,\qquad \K,\qquad \xi\F_p,
 \qquad \xi\F_{p^2}\quad(\xi\in\K^\times).
\end{equation}
The rank transition outside this set is
\begin{center}
\begin{tabular}{lccc}
\toprule
input & number & $\dim\J(A)$ & next state \\
\midrule
non-subfield plane, $p=2$ & $P-Q=30$ & $3$ & hyperplane \\
non-subfield plane, $p>2$ & $P-Q$ & $4$ & $\K$ \\
hyperplane & $L$ & $4$ & $\K$ \\
\bottomrule
\end{tabular}
\end{center}
On the recurrent scalar subfields,
$\J(\xi\F_{p^d})=\xi^{-1}\F_{p^d}$ for $d=1,2$; zero and $\K$ are fixed.

\item The maximum tail and depth enumerator
$D_p(u)=\sum_{A\in\X_p}u^{\tail(A)}$ are
\begin{equation}\label{eq:depth}
 \max_A\tail(A)=
 \begin{cases}2,&p=2,\\1,&p>2,\end{cases}
 \qquad
 D_p(u)=
 \begin{cases}
 R+Lu+(P-Q)u^2,&p=2,\\
 R+(S-R)u,&p>2.
 \end{cases}
\end{equation}
Moreover,
\begin{equation}\label{eq:image-stabilization}
 |\J^t(\X_p)|=
 \begin{cases}
 S,&t=0,\\
 R+L,&p=2,\ t=1,\\
 R,&p=2,\ t\geq2,\\
 R,&p>2,\ t\geq1.
 \end{cases}
\end{equation}

\item There are $R$ recurrent states and $F$ fixed states; the remaining
$R-F$ recurrent states form $(R-F)/2$ two-cycles.  Hence
\begin{equation}\label{eq:fix-iterates}
 |\Fix(\J^n)|=
 \begin{cases}F,&n\text{ odd},\\R,&n\text{ even},\end{cases}
 \qquad
 \zeta_{\J}(z)=(1-z)^{-F}(1-z^2)^{-(R-F)/2}.
\end{equation}
In particular, $F=4$ for $p=2$ and $F=6$ for odd $p$.

\item Every positive-time target fibre is
\begin{equation}\label{eq:fibres}
 \nu_t(B)=
 \begin{cases}
 1,&B\in\Rec(\J),\ B\ne\K,\\
 2,&p=2,\ t=1,\ \dim B=3,\\
 1+L,&p=2,\ t=1,\ B=\K,\\
 1+L+P-Q,&p=2,\ t\geq2,\ B=\K,\\
 1+L+P-Q,&p>2,\ B=\K,\\
 0,&\text{otherwise}.
 \end{cases}
\end{equation}

\item Every transient state lies in the component of the fixed state $\K$.
For odd $p$, that component is a star whose leaves are all hyperplanes and
non-subfield planes.  For $p=2$, its first level consists of the $15$
hyperplanes, each with exactly two non-subfield-plane children.  Every other
fixed point or two-cycle is a bare component.  Thus the number of weak
components is $(R+F)/2$.
\end{enumerate}
\end{theorem}

The residual theorem is the conjunction of the sharp characteristic
dichotomy, the complete graph and image stabilization, and the all-time
target fibres.  The recurrent classification, inverse-line geometry,
Gaussian counts, Singer action, and zeta conversion are treated as inputs or
standard consequences.  No novelty or priority inference is made.

\section{Rank growth and the inverse of a plane}

\begin{lemma}[rank monotonicity]\label{lem:rank}
For every $A\in\X_p$,
\begin{equation}\label{eq:rank}
 \dim\J(A)\geq\dim A.
\end{equation}
If equality holds, then
\begin{equation}\label{eq:equality}
 \J(A)=A^{-1}\cup\{0\},\qquad \J^2(A)=A.
\end{equation}
Consequently, $A$ is recurrent if and only if equality holds in
\eqref{eq:rank}, and every recurrent period divides two.
\end{lemma}

\begin{proof}
If $d=\dim A$, inversion gives $p^d-1$ distinct nonzero elements inside
$\J(A)$.  An $r$-dimensional subspace has only $p^r-1$ nonzero elements, so
$r\geq d$.  When $r=d$, cardinality forces the first identity in
\eqref{eq:equality}; patched inversion and a second span give the second.

Dimensions are nondecreasing along every orbit.  They must be constant on a
cycle, so recurrence forces equality at the first edge.  Conversely,
\eqref{eq:equality} places $A$ on a cycle of length at most two.
\end{proof}

We record the exact external input used with Lemma~\ref{lem:rank}.

\begin{proposition}[published inverse-image classification
\citep{KolomeecBykov2024}]\label{prop:external}
Let $E$ be a finite field of characteristic $p$, and patch inversion by
$0^{-1}=0$.  If an affine $\F_p$-subspace $A\subseteq E$ has $|A|>2$, then
its inverse image is an affine subspace if and only if
$A=\xi\F_{p^d}$ for a subfield $\F_{p^d}\subseteq E$ and
$\xi\in E^\times$.
\end{proposition}

Proposition~\ref{prop:external} receives zero contribution credit.  It
immediately rules out equality for a three-dimensional subspace of
$\F_{p^4}$, since a subfield degree must divide four.  The binary lines have
size two, outside the stated hypothesis, but their behavior is direct.

The next calculation exposes exactly where the prime two differs.  Its
projective interpretation is already contained in inverse-line geometry
\citep{FainaEtAl2002,LavrauwZanella2014}.

\begin{lemma}[plane span]\label{lem:plane}
Let $A$ be a plane in $\K$.  Then
\begin{equation}\label{eq:plane-rank}
 \dim\J(A)=
 \begin{cases}
 2,&A=\xi\F_{p^2},\\
 3,&A\ne\xi\F_{p^2}\text{ and }p=2,\\
 4,&A\ne\xi\F_{p^2}\text{ and }p>2.
 \end{cases}
\end{equation}
\end{lemma}

\begin{proof}
After scaling, write $A=\xi\langle1,\alpha\rangle_{\F_p}$ and put
$r=[\F_p(\alpha):\F_p]$.  Since $\alpha\in\F_{p^4}\setminus\F_p$, we have
$r\in\{2,4\}$.  Projective representatives for the inverse points, after
discarding the harmless scalar $\xi^{-1}$, are
\begin{equation}\label{eq:representatives}
 1,\qquad(\alpha-t)^{-1}\quad(t\in\F_p).
\end{equation}

Choose distinct $t_1,\ldots,t_s\in\F_p$ with $s+1\leq r$.  A relation
\[
 c_0+\sum_{i=1}^{s}\frac{c_i}{\alpha-t_i}=0
\]
becomes $G(\alpha)=0$ after multiplication by
$D(\alpha)=\prod_i(\alpha-t_i)$, where
\[
 G(x)=c_0D(x)+\sum_{i=1}^{s}c_i\frac{D(x)}{x-t_i}
\]
has degree at most $s<r$.  Minimality of $r$ gives $G=0$.  Evaluating this
polynomial identity at $t_i$ yields
$c_i\prod_{j\ne i}(t_i-t_j)=0$, so every $c_i$ and then $c_0$ vanish.
Thus the span of \eqref{eq:representatives} has dimension
$\min\{p+1,r\}$.

If $r=2$, then $\langle1,\alpha\rangle=\F_{p^2}$.  If $r=4$, the minimum is
three at $p=2$ and four for every odd prime.  Restoring $\xi$ proves
\eqref{eq:plane-rank}.
\end{proof}

\begin{corollary}[rank transitions]\label{cor:transitions}
Every hyperplane maps to $\K$.  Every non-subfield plane maps to a
hyperplane at $p=2$ and to $\K$ at odd $p$.
\end{corollary}

\begin{proof}
Lemma~\ref{lem:rank} leaves image dimension three or four for a hyperplane.
Equality at three contradicts Proposition~\ref{prop:external}, so its image
is $\K$.  Lemma~\ref{lem:plane} handles planes.
\end{proof}

\section{Sharp time and the recurrent core}

\begin{proof}[Proof of Theorem~\ref{thm:main}(i)--(iii), apart from the
image count at binary time one]
By Lemma~\ref{lem:rank}, a recurrent state of dimension at least two has a
patched inverse image that is a subspace.  Proposition~\ref{prop:external}
therefore makes it a scalar subfield.  The subfield degrees inside four are
$1,2,4$.  Lines, including binary lines, satisfy
$\J(\xi\F_p)=\xi^{-1}\F_p$ directly; the same identity holds for scalar
quadratic subfields.  This proves \eqref{eq:recurrent}.  Corollary
\ref{cor:transitions} supplies the remaining rows of the transition table.

There are $L$ lines, $P$ planes, and $L$ hyperplanes.  Scalar quadratic
subfields are parametrized by
$\K^\times/\F_{p^2}^\times$, so their number is
$(p^4-1)/(p^2-1)=Q$.  This proves the state count $S$ and recurrent count
$R$ in \eqref{eq:constants}.

Every hyperplane has tail one.  At an odd prime every non-subfield plane also
has tail one, whereas at $p=2$ it has tail two by Corollary
\ref{cor:transitions}.  Since
$P-Q=(p^2+1)(p^2+p)>0$, both bounds are attained.  Counting these strata
gives \eqref{eq:depth}; it also shows that $\J^2(\X_2)$ and
$\J(\X_p)$ for odd $p$ lie in the recurrent core.

On recurrent lines, $\J$ induces inversion on the cyclic quotient
$\K^\times/\F_p^\times$, which has order $L$.  On recurrent planes it
induces inversion on the cyclic quotient
$\K^\times/\F_{p^2}^\times$, of order $Q$.  Inversion on a cyclic group of
order $m$ has $\gcd(2,m)$ fixed points.  Adding zero and $\K$ gives $F$ fixed
states; Lemma~\ref{lem:rank} allows no periods beyond two.  Thus the other
$R-F$ recurrent states form $(R-F)/2$ two-cycles.  Since $L$ and $Q$ are odd
at $p=2$ and even at odd $p$, the displayed values of $F$ follow.

The fixed-iterate count in \eqref{eq:fix-iterates} is now immediate.  Taking
the exponential of its defining series, or multiplying the standard factors
for fixed points and two-cycles, gives the stated zeta function.
\end{proof}

The time jump is therefore a projective-cardinality effect.  It does not
come from a different rule in characteristic two: the literal update
\eqref{eq:map} is uniform in $p$.

\section{Every-target fibres and the component graph}

Only the binary plane-to-hyperplane row requires a target-resolved count.

\begin{lemma}[twisted scalar symmetry]\label{lem:singer}
For $\lambda\in\K^\times$ and $A\in\X_p$,
\begin{equation}\label{eq:twisted}
 \J(\lambda A)=\lambda^{-1}\J(A).
\end{equation}
The scalar action of $\K^\times$ is transitive on the hyperplanes of $\K$.
Consequently, at $p=2$ every hyperplane has exactly two non-subfield-plane
predecessors.
\end{lemma}

\begin{proof}
Equation \eqref{eq:twisted} follows by inverting every nonzero element before
taking the span.  The trace pairing on the separable extension
$\K/\F_p$ is nondegenerate.  Hence every hyperplane is
\[
 H_c=\{x\in\K:\Tr_{\K/\F_p}(cx)=0\}
\]
for some $c\in\K^\times$, unique modulo $\F_p^\times$.  Scalar
multiplication sends these kernels transitively to one another.

At $p=2$, Lemma~\ref{lem:plane} sends all $P-Q=30$ non-subfield planes to
hyperplanes.  Twisted scalar symmetry gives bijections between the fibres
over any two hyperplanes, so their sizes are constant.  There are $L=15$
targets, and $30/15=2$ sources lie over each.
\end{proof}

\begin{proof}[Completion of Theorem~\ref{thm:main}(ii), (iv), and (v)]
Lemma~\ref{lem:singer} shows that every binary hyperplane occurs at time one.
Together with the recurrent core this gives $|\J(\X_2)|=R+L$.  Every such
hyperplane maps to $\K$, so the second image is exactly the recurrent core.
At odd primes the transition table already sends every transient state to
$\K$, proving all of \eqref{eq:image-stabilization}.

The restriction of $\J$ to the recurrent core is an involutive bijection.
Every recurrent target other than $\K$ therefore has exactly one predecessor
at every positive time.  No transient state can be an additional predecessor:
the transition table sends every transient orbit toward $\K$.

A non-subfield plane has no predecessor.  Indeed, a source of smaller rank
cannot reach a plane, while a plane source with plane image would be an
equality case and hence recurrent.  At odd primes no update lands on a
hyperplane.  In the binary case Lemma~\ref{lem:singer} gives its two
predecessors at time one; one further update sends every one of those paths
to $\K$, so hyperplane fibres vanish for $t\geq2$.

At binary time one, the full field receives itself and all $L$ hyperplanes,
giving $1+L$ sources.  From time two onward it also receives all $P-Q$
non-subfield planes, giving $1+L+P-Q$.  At odd primes those planes arrive in
one step, so the latter count holds for every $t\geq1$.  This proves
\eqref{eq:fibres}.

The same argument identifies the components.  Every transient state reaches
$\K$, and \eqref{eq:fibres} shows that no transient tree attaches to any
other recurrent cycle.  For odd $p$ all transient vertices are leaves.  At
$p=2$, Lemma~\ref{lem:singer} gives the uniform two-level tree described in
part (v).  The recurrent core has $F$ fixed cycles and $(R-F)/2$ two-cycles,
so it has $(R+F)/2$ components.
\end{proof}

\section{Exact controls and scope}

The formulas specialize as follows.  The depth column lists the numbers at
successive tail depths, and the last column gives the full-field fibre at
times one and two.

\begin{center}
\begin{tabular}{c@{\quad}rrrrlrr}
\toprule
$p$ & $S$ & $|\im\J|$ & $R$ & $F$ & cycles & depths & $\nu_{1,2}(\K)$\\
\midrule
2 & 67 & 37 & 22 & 4 & $1^4\,2^9$ & $22,15,30$ & $16,46$\\
3 & 212 & 52 & 52 & 6 & $1^6\,2^{23}$ & $52,160$ & $161,161$\\
5 & 1120 & 184 & 184 & 6 & $1^6\,2^{89}$ & $184,936$ & $937,937$\\
\bottomrule
\end{tabular}
\end{center}

A paper-local Python verifier uses only the standard library.  It constructs
an irreducible quartic for each of $p=2,3,5$, implements field arithmetic,
enumerates every subspace in reduced row-echelon form, and recomputes every
directed edge.  It then checks the recurrent scalar-subfield classification,
rank transitions, cycles, depth histograms, images, twisted scalar symmetry,
and every target fibre at times one through four.  Two fresh processes match
the frozen $32{,}754$-assertion transcript byte for byte.  These exhaustive
finite checks are controls for implementation and boundary errors; the proofs
above establish the formulas.

The scope is deliberately narrow.  The base field is prime and the extension
degree is four.  The published inverse-image classification and inverse-line
geometry remain zero-credit inputs, and the cycle/zeta count is a short
consequence once the recurrent core is known.  What remains in the note is
only the exact temporal anomaly, complete graph and stabilization, and
all-time target-fibre synthesis for the literal map \eqref{eq:map}.  The
bounded source search gives no novelty or priority conclusion.  This anonymous
Round-0 artifact is \texttt{GREEN\_OWNER\_THIN / HOLD\_EXTERNAL}.

\bibliographystyle{plainnat}
\bibliography{references}

\end{document}
